\scanpage{92}
% Main text, lines 1--7; continuation from scan 91.
later on, between the respective possible notations
for the groups denoted here
$\underline M_{\rho,\sigma}$ and
$\underline M_{\rho,\sigma}/\underline\mu_\rho\simeq
\underline M_{\rho,\sigma}^{\sim}$
(and ditto with $\unsure{\natural|\sharp}$)
--- it is these latter, it is true, that seem to me the most interesting,
for it is they$^{*}$ that \unsure{link up with} the universal group
$\unsure{M|\mathcal M}=\unsure{N|\mathcal N}_{1,1}^{\sim}$
(extension of $\Gamma_{\mathbb Q}^{\sim}$ by
$\inserted{\unsure{\mathcal G|G}_{0}^{\wedge}=}
(\mathrm{SL}(2,\mathbb Z))^{\wedge}$).
\margin{rather, the group of their $\widehat{\mathbb Z}$-points
\unsure{respct...|respect...}}

% Main text, lines 8--25.
It would seem that these groups \unsure{would deserve} the simplest notation,
in this case $\underline M_{\rho,\sigma}$, and the groups denoted
previously $\underline M_{\rho,\sigma}$ would take on their $\sim$,
or an analogous sign. (The $\sim$ is already used in a rather
different sense in $\Gamma_{\mathbb Q}^{\sim}$,
$M_{0,3}^{\sim}$ etc --- on the other hand the superscripts
$',\deleted{\unsure{\natural|\sharp}},'',\,''',\,{}^{\circ},\,+,\,!,\,
\unsure{\natural|\sharp}$
\deleted{$\unsure{',\,',\,',\,',\,+}$}
are already used in different senses.
The double use of $\sim$ moreover risks leading to conflicts of
notation sooner or later, it is perhaps preferable to replace the $\sim$
in $\Gamma_{\mathbb Q}^{\sim}$ by letters such as
$\Gamma_{\mathbb Q}^{\text{ét}}$
(suggesting ``extended'') or \deleted{\unsure{gen}}
``generalized'' \unsure{gen\ldots}\ \deleted{\illegible{1}}
I will have to resolve these di-
\margin{91}

\scanpage{93}
\mspage{694}
% Main text, lines 1--5; continuation of di- on scan 92.
lemmas of notation, for the time being I keep the notations introduced for
\unsure{the need|the needs} in the present section, which will be
\unsure{revised|revised} with care in a special \S\ later on.

% Main text, lines 6--9.
We should still describe the subgroups
$\underline M_{\rho,\sigma}^{\sim}(-\bar j)$,
$\underline M_{\rho,\sigma}^{\prime\sim}(-\bar j)$
in $\underline M_{\rho,\sigma}^{\sim}$ written as a
$\frac12$ direct product, we find
\[
\begin{aligned}
(211)\qquad
\underline M_{\rho,\sigma}^{\sim}(-\bar j)
&\deleted{\unsure{\simeq\underline{\mathbb O}_{12}^{*}
(\mathbb G_a\times\underline\mu_G)}}
\inserted{\deleted{\unsure{\underline\mu(\lambda,x)}}}
\subset \underline M_{\rho,\sigma}^{\sim}
\simeq \mathbb O_{12}^{*}
(\mathbb G_a\times\underline{\unsure{\mathcal G|\widetilde G}}^{+})
\\[-1pt]
&=\left\{
\deleted{\underset{\mathbb O_{12}^{*}}
{\unsure{\underline\mu}}}
\left(
\underset{\substack{\rotatebox{90}{$\in$}\\\mathbb G_a}}{\lambda},
\underset{\substack{\rotatebox{90}{$\in$}\\
\underline{\unsure{\mathcal G|\widetilde G}}^{+}}}{x}
\right)\cdot
\underset{\substack{\rotatebox{90}{$\in$}\\
\mathbb O_{12}^{*}}}{\underline\mu}
\ \middle|\ x\cdot\deleted{\unsure{\underline\mu}}
\underbrace{u_{\underline\mu}\deleted{\unsure{\rho}}}
_{\in\underline G_{\rho,\sigma}}
\in N_\rho^{?}\right\}
\end{aligned}
\]
% Main text and calculation, lines 12--15.
where $u_{\underline\mu}$ is defined in (151), in particular
\[
u_{\underline\mu}(\rho)=\lambda_1^{\nu}
\unsure{\{\rho\}|\rho}\deleted{\unsure{=}}
\qquad\text{where}\quad \nu=\nu(\underline\mu),
\]
\[
\deleted{=\lambda_1^{3\nu''}\lambda_1^{-\nu_3}\rho}
\]
\deleted{hence}
\diagram{\begin{tikzcd}[ampersand replacement=\&,column sep=0.1em,row sep=1.6em]
\deleted{(x} \&
\deleted{u_{\underline\mu})(\rho)=x(\lambda_1)^{\unsure{\nu}}x(\rho)}
\\
\inserted{\unsure{\delta^{\nu''}\tau^{\nu_3}}}
\arrow[u,bend left=20] \&
\end{tikzcd}}
% Replacement calculation, lines 15--18.
we have seen (154)
\[
v=\underbrace{\sigma^{\nu_3}\delta^{-\nu''}u_{\underline\mu}}
_{\unsure{v|u}}\in\Gamma\underline N_\rho,
\qquad \underline\chi(v)=\underline\mu
\]
\margin{where $\delta=\lambda_1\lambda_0^{-1}$.}
and putting $\sigma=\omega\tau'\tau$, the first
\unsure{$\mathrm m^{\mathrm b}$|side} can also be written
\[
\deleted{(\unsure{\omega\tau'})^{\nu_3}\tau^{\nu_3}\unsure{u}}
\]
\[
(212)\qquad
v={\tau'}^{\nu_3}\left(\omega^{\nu_3}\cdot
\underbrace{\tau^{\nu_3}\delta^{-\nu''}u_{\underline\mu}}
_{\substack{
\in\Gamma\underline{\deleted{\unsure{N}}\inserted{G}}
_{\rho,\sigma}^{\unsure{\natural|\sharp}},\ \text{for the}\\
\text{three factors}\\
\text{are; hence}\\
\in\Gamma\underline N_\rho\cap
\underline G_{\rho,\sigma}^{\unsure{\natural|\sharp}}
=\Gamma\underline Z_\rho^{\unsure{\natural|\sharp}}
}}\right)
\]
\margin{92}

\scanpage{94}
% Main text, lines 1--4.
and we know that $v\in\deleted{\unsure{\underline Z}}
N_\rho^{?}$, hence the general form
\deleted{the \unsure{elements} of}
of the \deleted{elements}\inserted{sections}
of $N_\rho^{?}$ over
$\underline\mu=(\nu_2,\nu_3,\unsure{\gamma_0|\nu_0})
\in\Gamma\mathbb O_{12}^{*}$
\[
(212)\qquad
(\underline N_\rho^{?})_{\underline\mu}
=\left\{\omega^{i}\unsure{\sigma}^{\nu_3}\delta^{-\nu''}
u_{\underline\mu}\ \middle|\
i\in\Gamma\{0,1\}_{S_0}\right\}
\]
\unsure{with|where}
$\nu_3=\nu_3(\underline\mu)$,
$\nu''=\nu''(\underline\mu)$,
$\delta=\lambda_1\lambda_0^{-1}$. Thus we find
\[
\begin{aligned}
(213)\qquad
\underline M_{\rho,\sigma}^{\sim}(-\bar j)
=\left\{\left(
\underset{\substack{\rotatebox{90}{$\in$}\\\mathbb G_a}}{\lambda},
\underset{\substack{\rotatebox{90}{$\in$}\\
\underline{\unsure{\mathcal G|\widetilde G}}^{+}}}{x}
\right)\cdot
\underset{\substack{\rotatebox{90}{$\in$}\\
\mathbb O_{12}^{*}}}{\underline\mu}
\ \middle|\
\begin{gathered}
x=\omega^{i}\sigma^{\nu_3}\delta^{-\nu''}u_{\underline\mu},\ \text{where}\\
i\in\Gamma\{0,1\}_{S_0}
\end{gathered}\right\}
\end{aligned}
\]
\margin{$\delta=\lambda_1\lambda_0^{-1}$}

We will therefore have
\[
\begin{aligned}
(214)\qquad
\underline M_{\rho,\sigma}^{\prime\sim}(-\bar j)
&=\underline M_{\rho,\sigma}^{\sim}(-\bar j)\cap
\underline M_{\rho,\sigma}^{\unsure{\natural|\sharp}}\\
&=\left\{(\lambda,\delta^{-\nu''})\underline\mu
\ \middle|\
\begin{gathered}
\underline\mu\in\Gamma\mathbb O_{12}^{*\prime}
\quad(\text{hence }\varepsilon_3(\underline\mu)=\unsure{0|1})\\
\lambda\in\Gamma\mathbb G_a\\
\nu''=\nu''(\underline\mu)
\end{gathered}\right\}
\end{aligned}
\]
\margin{93}

\scanpage{95}
\mspage{695}
% Heading and opening, lines 1--3.
$11^{\mathrm o})$ The homomorphism
$M\longrightarrow\underline M_{\rho,\sigma}(\widehat{\mathbb Z})^{\Delta}$,
relation with the tower of Fermat curves.

Let us note that we have for every integer $N\in\mathbb N^{*}$
\[
(215)\qquad
\mathbb O_{N}^{\Delta}
(\widehat{\mathbb Z})\simeq\widehat{\mathbb Z}
\qquad\text{hence}\qquad
\mathbb O_{N}^{*}(\widehat{\mathbb Z})^{\Delta}
\simeq\widehat{\mathbb Z}^{*}.
\]
\margin{The $\Delta$ indicates that we take the sections which are constant
modulo the identity component.}

We therefore have, taking into account the decompositions into semidirect products (150) etc,
isomorphisms
\[
(216)\qquad
\left\{\begin{aligned}
Z_{\rho,\sigma}^{\unsure{\ell|\wedge}}
&\overset{\mathrm{dfn}}{=}
\underline Z_{\rho,\sigma}^{\deleted{\unsure{\natural|\sharp}}}
(\widehat{\mathbb Z})^{\Delta}
=\widehat{\mathbb Z}^{*}\cdot
\deleted{\unsure{\widehat{\mathbb Z}}}\widehat L_0
\quad\text{where}\quad
\widehat L_0\simeq\widehat{\mathbb Z}\cdot\lambda_0\simeq\widehat{\mathbb Z}
\\
G_{\rho,\sigma}
&\overset{\mathrm{dfn}}{=}
\underline G_{\rho,\sigma}^{\deleted{\unsure{\Delta|\natural}}}
(\widehat{\mathbb Z})^{\Delta}
=\widehat{\mathbb Z}^{*}\cdot
\underline{\unsure{\mathcal G|\widetilde G}}^{+\wedge}
\\
G_{\rho,\sigma}^{\unsure{\natural|\sharp}}
&\overset{\mathrm{dfn}}{=}
\underline G_{\rho,\sigma}^{\unsure{\natural|\sharp}}
(\widehat{\mathbb Z})^{\Delta}
\simeq\widehat{\mathbb Z}^{*}\cdot\widehat M
\\
N_{\rho\deleted{\unsure{,\sigma}}}^{\sim\deleted{\unsure{\natural}}}
&\overset{\mathrm{df}}{=}
\underline N_\rho^{\sim}(\widehat{\mathbb Z})^{\Delta}
\simeq\mathbb O_{12}^{*}
(\widehat{\mathbb Z})^{\Delta}\simeq\widehat{\mathbb Z}^{*}
\\
N_\sigma^{\unsure{\natural|\sharp}}
&\overset{\mathrm{df}}{=}
\underline N_\sigma^{\unsure{\natural|\sharp}}
(\widehat{\mathbb Z})^{\Delta}
\simeq\deleted{\mathbb O_{12}^{*}}
\widehat{\mathbb Z}^{*}\cdot\widehat{\mathbb Z}
\\
\underset{\deleted{\unsure{\infty}}}{M}_{\rho,\sigma}^{\sim}
&\overset{\mathrm{df}}{=}
\underline M_{\rho,\sigma}^{\deleted{\unsure{\natural|\sharp}}}
(\widehat{\mathbb Z})^{\Delta}
\simeq\widehat{\mathbb Z}^{*}\cdot
(\unsure{\widehat{\mathbb Z}}\times
\unsure{\mathcal G|\widetilde G}^{+\wedge})
\\
M_{\rho,\sigma}^{\sim\unsure{\natural|\sharp}}
&\overset{\mathrm{df}}{=}
\underline M_{\rho,\sigma}^{\unsure{\natural|\sharp}}
(\widehat{\mathbb Z})^{\Delta}
\simeq\widehat{\mathbb Z}^{*}\cdot
(\widehat{\mathbb Z}\times\widehat M)
\end{aligned}\right.
\]
% Main text, lines below (216).
where the operations of $\widehat{\mathbb Z}^{*}$ on the
\unsure{$2^{\text{nd}}$|$2^{\text{nd}}$}
$\frac12$ direct factors are made explicit in the preceding sections.
We will note that, despite the \unsure{trappings|appearances} of the
notations and definitions, the \inserted{profinite} groups
$Z_{\rho,\sigma}$ and preceding ones are also those which would be
\unsure{constructed} \deleted{\unsure{$\mathcal G$}}
directly, in terms of
$(\unsure{\mathcal G|\widetilde G}^{\unsure{+|\natural}\wedge},
\unsure{\mathcal G|\widetilde G}^{\unsure{\natural|\sharp}\wedge}
=M^{\wedge})$
(equipped$^{*}$ with $\rho,\sigma,\tau$), in \S\ 49 $\uline{\overline{\mathrm{VI}}}$.
We therefore have a canonical homomorphism
\diagram{\[
(217)\qquad
\left\{\begin{gathered}
\begin{tikzcd}[ampersand replacement=\&,column sep=5em,row sep=2.5em]
M=\unsure{N|\mathcal N}_{1,1}^{\text{ét}}
\arrow[r]\arrow[d,phantom,"\cup" description]
\& M_{\rho,\sigma}^{\sim}\arrow[d,phantom,"\cup" description]
\\
M^{!}\arrow[r] \& M_{\rho,\sigma}^{\sim\unsure{\natural|\sharp}}
\end{tikzcd}
\end{gathered}\right.
\]}
which as usual decomposes into
\margin{94}